\documentclass[11pt]{article}
\usepackage[a4paper,margin=23mm]{geometry}
\usepackage{amsmath,amssymb,amsthm,booktabs,array}
\usepackage[hidelinks]{hyperref}
\newtheorem{theorem}{Theorem}
\newtheorem{lemma}{Lemma}
\newcommand{\SF}{\operatorname{SF}_2}
\newcommand{\bb}{\beta_2}
\setlength{\parindent}{0pt}
\setlength{\parskip}{5pt}
\begin{document}
\begin{center}
{\Large Exact star formation under arbitrary interfaces}\\[4pt]
Research note for discussion, 15 October 2026
\end{center}

For a finite simple graph, a set $S$ is \emph{2-star-forming} if every
$x\notin S$ lies in a copy of $K_{1,2}$ in $G[S\cup\{x\}]$. Copies need not
be induced. Write $\SF(G)$ for the largest size of an inclusion-minimal such
set, and $\bb(G)$ for the largest size of a set inducing maximum degree at
most one. If $d_S(x)=|N(x)\cap S|$, the exact local star test is
\[
 d_S(x)\ge2\quad\text{or}\quad
 \exists u\in N(x)\cap S:\ d_S(u)\ge1. \tag{1}
\]
The alternatives make $x$ the centre or a leaf. Star formation is upward
closed, so inclusion-minimality is equivalent to failure of all single
deletions. The tested failure vertex may differ from the deleted vertex.

\textbf{The fixed graph.} Let $H$ have parts $A=\{a_1,\ldots,a_5\}$ and
$B=\{b_1,\ldots,b_5\}$, and cross-edge matrix
\[
\begin{pmatrix}
1&1&0&0&1\\1&1&0&1&0\\0&0&0&1&1\\0&1&1&1&1\\1&0&1&1&1
\end{pmatrix}.
\]
Degrees on both sides are $(3,3,2,4,4)$. Put
$S_0=\{a_1,a_2,a_3,b_1,b_2,b_3\}$. Each outside vertex has two neighbours
in $S_0$. The following single deletions fail (1).
\begin{center}
\begin{tabular}{@{}cccccc@{}}
\toprule
Deleted & $a_1$ & $a_2$ & $a_3$ & $b_1$ & $b_2$ \\
Failure & $b_5$ & $b_4$ & $a_3$ & $a_5$ & $a_4$ \\
Remaining neighbours & $\{a_3\}$ & $\{a_3\}$ & $\varnothing$ & $\{b_3\}$ & $\{b_3\}$ \\
\bottomrule
\end{tabular}
\end{center}
Deleting $b_3$ fails at $b_3$, with no remaining selected neighbour.
Vertices $a_3,b_3$ are isolated in $H[S_0]$, so the singleton neighbours
in the table cannot act as centres. Thus $S_0$ is minimal and $\SF(H)\ge6$.

\textbf{Dissociation bound.} The matrix is symmetric under exchanging
$a_i,b_i$. A dissociation six-set has side split $(1,5)$, $(2,4)$ or $(3,3)$,
up to this symmetry. The first split violates minimum degree two. In
$(2,4)$, omitting one $B$ vertex can lower selected degree to at most one
only for $a_3$, so two selected $A$ vertices are impossible. In $(3,3)$,
the degree-four vertices $a_4,a_5$ are excluded, forcing $a_1,a_2,a_3$.
Both $a_1,a_2$ then need their common pair $b_1,b_2$ omitted, leaving both
$b_4,b_5$ adjacent to $a_3$. Thus $\bb(H)\le5$; the independent set $A$
attains five.

\begin{theorem}
Let $I$ be any finite set and $R\subseteq I\times I$. Take disjoint copies
$H_i$ and add $b_1^{(i)}a_1^{(j)}$ for every $(i,j)\in R$, obtaining $F_R$.
Then $\bb(F_R)=5|I|$ and $\SF(F_R)=6|I|$.
For a directed cycle on $r\ge2$ indices, $F_R$ is two-connected and
bipartite, has maximum degree four, order $10r$, and gap exactly $r$.
\end{theorem}
The upper bound, including $\SF(H)\le6$, follows from the local argument
on the next two pages. No enumeration is needed for these assertions.

\newpage
\textbf{A dense selection admits a safe internal deletion.}
Only the ports $p=a_1,q=b_1$ may meet vertices outside a copy of $H$.

\begin{lemma}
For $T\subseteq V(H)$ with $|T|\ge6$, star formation is equivalent to
domination: every vertex outside $T$ has a neighbour in $T$.
\end{lemma}
\begin{proof}
Star formation implies domination. Conversely, suppose $x\notin T$ has
the unique selected neighbour $u$, and $u$ has no neighbour in $T$. Then
both $N(x)\setminus\{u\}$ and $N(u)$ lie in $Q=V(H)\setminus T$.
They are disjoint by bipartiteness. Every edge $xu$ satisfies
$d(x)+d(u)\ge6$: the degree-two vertices have only degree-four neighbours,
and other endpoints have degree at least three. Consequently
\[
 |Q|\ge d(x)-1+d(u)\ge5,
\]
contradicting $|Q|\le4$. Thus every singleton neighbour can serve as a
centre; vertices with two selected neighbours are themselves centres.
\end{proof}

The closed neighbourhoods of size at most four are precisely
\begin{align*}
 Z_{a_3}&=\{a_3,b_4,b_5\},& Z_{b_3}&=\{b_3,a_4,a_5\},\\
 Z_{a_1}&=\{a_1,b_1,b_2,b_5\},& Z_{a_2}&=\{a_2,b_1,b_2,b_4\},\\
 Z_{b_1}&=\{b_1,a_1,a_2,a_5\},& Z_{b_2}&=\{b_2,a_1,a_2,a_4\}.
\end{align*}
The two triples are disjoint, and any two distinct displayed quadruples
intersect in at most two vertices. All other closed neighbourhoods have
size five.

\begin{lemma}
If $U\subseteq V(H)$ dominates $H$ and $|U|\ge7$, there is
$v\in U\setminus\{p,q\}$ such that $U\setminus\{v\}$ is star-forming
inside $H$, and each selected port having a selected internal neighbour
before deletion still has one afterwards.
\end{lemma}
\begin{proof}
Put $W=V(H)\setminus U$, so $|W|\le3$. A candidate deletion $v$ destroys
domination exactly when some $N[x]$ is contained in $W\cup\{v\}$.
Since $U$ dominates, $v$ must be the unique element of $N[x]\setminus W$.

At most one candidate completes one of the two triples: each requires two
members of that triple already in $W$, whereas the triples are disjoint.
At most one candidate completes a quadruple: this requires $|W|=3$ and
$W$ contained in that quadruple, and two distinct quadruples cannot both
contain the same three-element $W$. A size-five neighbourhood is too
large. At most two candidates are therefore bad for domination.

For each port, at most one further candidate destroys its internal
selected-neighbour activity, namely its unique neighbour in $U$, if it
has one. Thus at most four candidates are bad in total. There are at least
$|U|-2\ge5$ candidates outside the ports, so a safe candidate remains.
Its deletion preserves domination and leaves at least six vertices;
Lemma 1 gives star formation. It also preserves port activity as required.
\end{proof}

The count is a union bound on bad deletion choices. Overlaps are harmless;
no surplus or witness is assigned twice to justify different deletions.

\newpage
\textbf{The local bound in an arbitrary surrounding graph.}

\begin{lemma}
If $G$ contains an induced $H$ whose external edges touch only $p,q$, then
every inclusion-minimal 2-star-forming set $S$ of $G$ satisfies
$|S\cap V(H)|\le6$.
\end{lemma}
\begin{proof}
Suppose $U=S\cap V(H)$ has size at least seven. First $U$ dominates $H$.
Otherwise some $x\notin U$ has $N[x]\subseteq V(H)\setminus U$, a set of
size at most three. The list of closed neighbourhoods forces $x=a_3$
or $b_3$. These are not ports, so $x$ has no selected neighbour anywhere
in $G$, contradicting star formation.

Choose $v$ from Lemma 2 and set $S'=S\setminus\{v\}$. Every vertex of
the copy outside $S'$ has a star entirely inside the copy. For a vertex
$x$ outside the copy and outside $S'$, take its old witnessing star.
If that star avoids $v$, it survives. Otherwise $x$ is not adjacent to the
non-port $v$, so $v$ must be a leaf, the centre must be a selected port
$u$, and the other leaf must be $x$. The port had the selected internal
neighbour $v$, so after deletion it has some selected internal neighbour
$w$. The star with centre $u$ and leaves $x,w$ survives. Thus $S'$ is
star-forming, contradicting minimality.
\end{proof}

\textbf{Proof of the theorem.} Applying Lemma 3 to each copy bounds any
minimal set by $6|I|$. The union of the sets $S_0$ attains this size and
remains minimal: the outside vertices keep their two selected neighbours,
and all six deletion witnesses from page 1 persist. Interfaces only meet
$a_1,b_1$, so the possible singleton centres $a_3,b_3$ remain isolated.
Upward closure then proves minimality. This gives $\SF(F_R)=6|I|$.

Every dissociation set restricts to at most five vertices in each induced
copy. The full $A$ part attains $5|I|$, since every interface joins $B$ to
$A$. Hence $\bb(F_R)=5|I|$. Loops of $R$ only specify the already existing
edge $b_1a_1$ and do not affect the argument.

Each copy contains the Hamilton cycle
\[
 a_1,b_1,a_2,b_2,a_4,b_3,a_5,b_4,a_3,b_5,a_1.
\]
In the cyclic family, deleting a vertex leaves its copy connected and
destroys at most one interface, because the two ports are distinct.
The remaining interfaces still connect all copies. For $r=2$ at least one
of the two distinct interfaces remains. Thus the graph is two-connected.
The ports rise from degree three to four and all other degrees stay fixed.
\hfill$\square$

\textbf{Verification and scope.} Independent base enumeration returns
$\bb(H)=5$, $\SF(H)=6$ and 37, 48, 1 minimal sets of sizes 4, 5, 6.
Separate controls check all 174 dominating selections of size at least
seven and 29,656 safe global deletions in 60 surrounding graphs. These
finite controls check the argument; the all-size conclusion follows from
the proofs above. The accompanying Lean sources check the literal/local
definition, witness preservation and the general local-deletion argument;
the precise completed declarations are listed with the source release.
This note does not settle a sharp universal gap coefficient, computational
complexity, or the rectangular-grid triple equality.
\end{document}
